\chapter{Nonlinear equations}\label{ch:nonlin}
\chapterbrief
{a scalar evolution in one space variable, $u_t=\sum_nT_n(u)$, each term $T_n$ linear or polynomial in the jet $(u,u_x,u_{xx},\dots)$; a smooth datum $u_0$; the exact solution is known. Benchmarks: viscous Burgers, KdV--Burgers, Fisher--KPP; KdV and NLS as integrable controls.}
{for a nonlinear equation the single-term flows do not commute and superposition fails. ``Component per term'' needs a new definition, and identification needs a new frame (a dictionary). We give both and compute them where everything is explicit.}
{fourth link of the chain. It reuses the earlier chapters through the jet expansion and the annihilator idea. Its results are weaker and relative, but sharp on the benchmark models.}

\section{The interaction term for nonlinear flows}

Let $S_n^t$ be the flow of $u_t=T_n(u)$ alone, which is now a nonlinear map. Definition~\eqref{eq:eint-def} still makes sense:
\[
\Eint(t)=u(t)-u_0-\sum_n\big(S_n^tu_0-u_0\big).
\]
There is no product formula, because the flows do not commute. But the Taylor expansion in time is explicit.

\begin{theorem}{Jet of the nonlinear interaction \led{S4T-T1} \proved}\label{thm:jet}
For smooth data,
\[
\Eint(t)=\frac{t^2}{2}\sum_{n\ne m}T_n'[u_0]\big(T_m(u_0)\big)+O(t^3),
\]
where $T_n'[u_0](v)$ is the derivative (Fr\'echet derivative) of $T_n$ at $u_0$ in the direction $v$.

\dimref{S4T-T1}
\end{theorem}

\begin{intuition}
To first order in $t$ every flow moves $u_0$ by $t\,T_n(u_0)$, and these moves add exactly: there is no interaction at order $t$. At order $t^2$ the full flow sees each term react to the motion caused by the others. That reaction is the derivative of $T_n$ in the direction of $T_m(u_0)$.
\end{intuition}

\begin{example}{Burgers}\label{ex:burgers-jet}
For $u_t=-uu_x+\nu u_{xx}$ take $T_1(u)=-uu_x$ and $T_2(u)=\nu u_{xx}$. The derivatives are
\[
T_1'[u](v)=-(vu_x+uv_x),\qquad T_2'[u](v)=\nu v_{xx}.
\]
Substituting the other term:
\begin{align*}
T_1'[u_0]\big(T_2(u_0)\big)&=-\nu\,\big(u_0''u_0'+u_0u_0'''\big),\\
T_2'[u_0]\big(T_1(u_0)\big)&=-\nu\,\big(3u_0'u_0''+u_0u_0'''\big).
\end{align*}
Half the sum, times $t^2$, gives
\[
\Eint(t)=-\nu\,\big(u_0u_0'''+2u_0'u_0''\big)\,t^2+O(t^3).
\]
\end{example}

\section{Before, at and after the shock}

The flow of the single term $-uu_x$ (inviscid Burgers) forms shocks: along each characteristic $u$ is constant and the characteristic moves with speed $u$, so faster parts overtake slower ones. For $u_0=\sin x$ the characteristics are
\[
x=\xi+t\sin\xi,
\]
and since
\[
\frac{\partial x}{\partial\xi}=1+t\cos\xi
\]
first vanishes at $\xi=\pi$, the breaking time is $t^*=1$ (Figure~\ref{fig:burgers}).

\begin{figure}[!htb]
\centering
\begin{tikzpicture}
\begin{axis}[width=9.8cm,height=5.2cm,xlabel={$x$},ylabel={$u$},xmin=0,xmax=6.283,ymin=-1.9,ymax=1.9,
 legend style={font=\scriptsize,draw=none,fill=white,fill opacity=0.85,text opacity=1,at={(0.01,0.02)},anchor=south west},
 title={\small Inviscid flow of $-uu_x$, $u_0=\sin x$},no markers,domain=0:6.283185,samples=200,smooth]
\addplot[cG,thick,dashed,variable=\z] ({\z},{sin(deg(\z))});
\addplot[cA,thick,variable=\z] ({\z+0.5*sin(deg(\z))},{sin(deg(\z))});
\addplot[cC,thick,variable=\z] ({\z+1.0*sin(deg(\z))},{sin(deg(\z))});
\addplot[cB,thick,variable=\z] ({\z+1.6*sin(deg(\z))},{sin(deg(\z))});
\draw[cB,dashed] (axis cs:3.14159,-1) -- (axis cs:3.14159,1);
\legend{$t=0$,$t=0.5$,$t=1$ (breaking),$t=1.6$ (multivalued)}
\end{axis}
\end{tikzpicture}
\caption{At $t=1$ the profile has a vertical tangent and the gradient blows up. Beyond that the curve folds; the physical (entropy) solution replaces the fold by a jump at $x=\pi$. This breaking is what makes the interaction term discontinuous after the shock.}\label{fig:burgers}
\end{figure}

The interaction changes qualitatively with time, and each regime needs different tools.

\subsection{Before the shock: a closed form}

\begin{theorem}{Closed form before the shock \led{S4T-T2, S4T-P3} \proved}\label{thm:preshock}
For Burgers with $t<t^*$, the interaction term has a closed form at first order in the viscosity, with remainder
\[
O\big(\nu^2\varepsilon^2t^3\big),
\]
where $\varepsilon$ measures the size of the datum. The closed form blows up at the breaking time $t^*$.

\dimref{S4T-T2}
\end{theorem}

The tool is the Cole--Hopf substitution
\[
u=-2\nu\,\frac{\varphi_x}{\varphi},
\]
which turns Burgers into the heat equation $\varphi_t=\nu\varphi_{xx}$~\cite{ch}. It plays the role of an ``exact linearisation''. For KdV and NLS the inverse scattering transform plays the same role.

The analyticity radius of the Burgers solution with datum $U_0$ is at least $2/\mathrm{osc}(U_0)$, and this bound is optimal \led{F4L-P3} \proved.

\subsection{After the shock: sizes in $L^1$}

The flow of $-uu_x$ alone forms a shock, so $\Eint$ has jumps whose size does not depend on $\nu$ \led{S4T-T3} \modulo{[F1]}. This is forced by the definition. A function with a jump is not in $H^1$, so Sobolev norms are no longer the right measure. Pointwise values are unreliable, but the \emph{measure} of the discrepancy is well behaved. The meaningful sizes are in $L^1$.

\begin{theorem}{Interaction after the shock \led{OPA-T1} \modulo{Kru\v zkov}}\label{thm:shock}
At a fixed time $t>t^*$, under a nondegeneracy condition $(H_t)$ that holds for example for $u_0=\sin x$,
\[
\|\Eint(t)\|_{L^1}=\Theta(\nu)\qquad(\nu\to0).
\]
Each shock of strength $\rho$ contributes
\[
4\nu\log\big(2\cosh(\rho/2)\big)
\]
to the leading term.

\dimref{OPA-T1}
\end{theorem}

\begin{theorem}{Uniform rate across the shock time \led{OPF5A-T1, OPF5A-T2} \modulo{Kru\v zkov}}\label{thm:uniform}
Uniformly in time, near the shock time,
\[
\|\Eint(t)\|_{L^1}\approx2\nu\log\frac{1}{\max(|t-t^*|,\nu)}+O(\nu),
\]
so that
\[
\sup_{t\le T}\|\Eint(t)\|_{L^1}=\nu|\log\nu|+O(\nu),
\]
with leading coefficient exactly $1$. The constant at the shock time is explicit (it involves the Pearcey constant $1.188654\ldots$). For data with several shock births the uniform rate is \condl{} on a localisation lemma.

\dimref{OPF5A-T2}
\end{theorem}

\begin{proposition}{The splitting error is larger \led{OPF5D-T6} \proved}\label{prop:splitting}
The error of the splitting scheme ``advection, then diffusion'' is
\[
\|S(t)u_0-S_2^tS_1^tu_0\|_{L^1}=\Theta\big(\sqrt\nu\big),
\]
larger than the interaction term itself. Treating the terms one after the other costs more than ignoring their interplay.

\dimref{OPF5D-T6}
\end{proposition}

\begin{table}[!htb]
\centering\small
\begin{tabular}{@{}L{3.0cm}L{4.8cm}L{6.4cm}@{}}
\toprule
Regime & What happens & Tools and results\\\midrule
before, $t<t^*$ & smooth; characteristics do not cross & jet (Thm.~\ref{thm:jet}), closed form (Thm.~\ref{thm:preshock}), Cole--Hopf\\
at breaking, $t=t^*$ & the gradient blows up & smoothness estimates fail; the closed form blows up\\
after, $t>t^*$ & the single-term flow has a jump & $L^1$: $\Theta(\nu)$, uniform $\nu|\log\nu|$, splitting $\Theta(\sqrt\nu)$\\
\bottomrule
\end{tabular}
\caption{The three regimes of the Burgers interaction term.}\label{tab:regimes}
\end{table}

\section{Identification relative to a dictionary}

Nonlinear terms cannot be handled by the ideal $\Ann(u)$: there is no ring of operators. Instead one fixes a finite list of candidates.

\begin{notebox}{Dictionary setting}
\begin{itemize}
\item A \emph{dictionary} $\mathcal D$ is a finite set of monomials in the jet $(u,u_x,u_{xx},\dots)$ of degree at most $d$ and order at most $k$, for example $\{u,\ u_x,\ uu_x,\ u_{xx},\ u_{xxx}\}$. $V=\spn\mathcal D$.
\item $I(u)$ is the set of all differential polynomial relations that $u$ satisfies.
\item $[F]$ is the differential ideal generated by the equation $F$: $F$, its derivatives $\partial_xF,\partial_tF,\dots$, and their combinations.
\item $\mathcal J_u$ is a jet-level obstruction defined from $u$, and $\psi$ maps the dictionary into the jet space.
\end{itemize}
\end{notebox}

\begin{theorem}{Dictionary criterion \led{S4ID-T1} \proved}\label{thm:dict}
For smooth $u$ on any $\Omega$ there is an exact sequence
\[
0\to[F]\cap V\to I(u)\cap V\to\mathcal J_u\cap\psi(V)\to0 .
\]
In particular
\[
I(u)\cap V=[F]\cap V\iff\mathcal J_u\cap\psi(V)=0 .
\]
\dimref{S4ID-T1}
\end{theorem}

The correct statement is therefore ``the relations are exactly $F$ \emph{modulo its differential consequences}''. The literal statement $I(u)\cap V=\spn\{F\}$ is false in complete dictionaries, because they contain $\partial_xF$ \led{S4ID-P1} \proved.

\begin{theorem}{Burgers with $N$ shocks \led{S4ID-T2, OP-T7} \proved}\label{thm:burgersN}
For the $N$-shock solutions of Burgers, $u$ identifies its equation in the evolutionary dictionary with $d$ unknown terms if and only if $d\le N$.

\dimref{S4ID-T2}
\end{theorem}

\begin{example}{One shock, two unknowns}
The travelling front $u=-\tanh(x/(2\nu))$ is a $1$-shock solution. It satisfies Burgers, but also the ODE that every travelling wave satisfies. With $d=2$ unknown coefficients (in front of $uu_x$ and $u_{xx}$) and $N=1$, the theorem says $u$ does not identify the pair. It fixes only their ratio.
\end{example}

\begin{theorem}{KdV $N$-solitons \led{S4ID-T3, OPF5C-T2} \modulo{Novikov 1974}}\label{thm:kdv}
For the $N$-soliton solution of KdV, the annihilator in the jet variables is principal, generated by one polynomial $\Pi_N$ of total degree
\[
\deg\Pi_N=(N+2)\,2^{N-1}.
\]
For generic $N=2$ this degree is $8$. The solution identifies KdV in a dictionary of degree $d$ if and only if $d<(N+2)2^{N-1}$.

\dimref{OPF5C-T2}
\end{theorem}

\begin{proposition}{A closed form does not identify \led{S4ID-T4} \proved}
If $u$ has a closed form, it lives on a smaller variety and satisfies extra relations: it does not identify its equation in a complete dictionary.

\dimref{S4ID-T4}
\end{proposition}

This is the nonlinear analogue of ``exponential-polynomials never identify''.

\section{Integrable models and scattering}

\begin{itemize}
\item \textbf{Hirota $N$-solitons.} The interaction is pairwise and localised: it is determined by pairs of solitons, integrates to $0$, and costs $O(N^3)$ \led{INT-P5, INT-P6} \proved.
\item \textbf{Scattering identification.} For reflectionless KdV, Novikov's centraliser characterises the potential \led{INT-T1, OPF5D} \proved. For NLS with multiplicity, identification by the scattering data holds; without multiplicity it is \falsest. The Lax and zero-curvature relations hold for all $n$ \led{OP-P17, OP-P18} \proved.
\item \textbf{KdV--Burgers.} The jet is proved. Before the shock the dispersive correction has remainder $O\big((\nu+|\delta|)^2\big)$ \led{OPR2-T5} \modulo{local well-posedness}.
\item \textbf{Fisher--KPP.} Only a jet-level statement with an estimate that is not uniform in time \led{S4T-P5} \proved. No sharp analogue of the Burgers asymptotics is claimed.
\end{itemize}

\begin{caution}[title={Open in the nonlinear part}]
For analytic KdV data with nonzero reflection, is there an autonomous ODE in $x$? Three obstructions are proved, among them the cubic relation of Airy \openq{} \led{INT-Q1}. The exact minimal syzygy degrees $\delta_N(K)$ are known in nine computer-assisted cases; their general formulas are conjectures \openq{} \led{OPF5C-Q1}.
\end{caution}

\section{Linear versus nonlinear}

\begin{table}[!htb]
\centering\small
\begin{tabular}{@{}L{3.0cm}L{5.8cm}L{5.8cm}@{}}
\toprule
 & Linear (Chapters~\ref{ch:algebra}--\ref{ch:interaction}) & Nonlinear (this chapter)\\\midrule
Candidate equations & the ring $\ring=\C[\tau,\xi]$ & a finite dictionary $V=\spn\mathcal D$\\
Relations of $u$ & the ideal $\Ann(u)$ & the differential ideal $I(u)$\\
What $P$ generates & the principal ideal $(P)$ & the differential ideal $[F]$\\
Identification & $\Ann(u)=(P)$ & $I(u)\cap V=[F]\cap V$\\
What escapes & the point spectrum $\Var(J_u)$ & the obstruction $\mathcal J_u\cap\psi(V)$\\
Solutions that never identify & exponential-polynomials & closed forms\\
Interaction & product formula, $O(N)$ & $t^2$ jet; $\Theta(\nu)$, $\nu|\log\nu|$ after shocks\\
\bottomrule
\end{tabular}
\caption{The linear and nonlinear problems side by side.}\label{tab:lin-nonlin}
\end{table}

\begin{keypoint}
Nonlinearity changes the objects (single-term flows with shocks, dictionaries instead of rings) but not the question. Where the formula is closed (before the shock, solitons) the answer is explicit. Across a shock the interaction does not vanish with the viscosity: it shrinks like $\nu|\log\nu|$.
\end{keypoint}
